\documentclass[11pt,a4paper]{article}
\usepackage[T1]{fontenc}
\usepackage{lmodern}
\usepackage[margin=2.5cm]{geometry}
\usepackage{amsmath}
\usepackage{graphicx}
\usepackage{authblk}
\usepackage[numbers,sort&compress]{natbib}
\usepackage[hidelinks]{hyperref}

\title{Title of the Short Communication}
\author[1]{First Author}
\author[2]{Second Author}
\affil[1]{Department, University, City, Country}
\affil[2]{Institute, City, Country}
\date{}

\begin{document}

\maketitle

\begin{abstract}
One paragraph of up to 100 words: the finding and why it matters.
\end{abstract}

\noindent\textbf{Keywords:} first keyword, second keyword

\section*{Main text}
A short communication reports one new finding in about 1500 words. State the
question~\cite{shannon1948}, what was done and what was found
(Figure~\ref{fig:overview}).

\begin{figure}[htbp]
  \centering
  \includegraphics[width=0.6\linewidth]{example-image}
  \caption{Overview of the method. Upload an image and use its file name here.}
  \label{fig:overview}
\end{figure}

State what the finding means and what should be done next.

\section*{Data availability}
Say where the data and code can be found.

\bibliographystyle{unsrtnat}
\bibliography{references}

\end{document}
